\section{Limitations and future work}

\subsection{Limitations}

The construction has several serious limitations.

\paragraph{Depth growth.}
Row support grows as $w_D=w_0^{2^D}$.  This restricts \sable{} to low-depth computation unless a new compression or bootstrapping method is added.

\paragraph{Low-noise tension.}
Correctness prefers small $\eta$ and $\eta_c$, but security against low-noise LPN attacks may prefer larger noise or larger dimensions.  This tension is central and must be quantified.

\paragraph{Large evaluation keys.}
The expansion key contains $N$ GSW matrices, each with $N$ sparse rows.  The compaction key also contains $N$ code/LPN ciphertexts.  Seeded randomness and structured sampling may reduce storage, but those optimizations need separate security analysis.

\paragraph{Secret-key setting.}
This draft gives a secret-key HE scheme with public evaluation keys.  Public-key encryption is not included.  A public-key variant may be possible using standard LPN public-key techniques, but it would need a new proof.

\paragraph{Compactor quality.}
A repetition-code compactor is simple but inefficient.  A serious version should use stronger q-ary codes and possibly CRT decomposition, inspired by recent code-based secure aggregation work \cite{BitzerEggerLiuWachterZeh2026}.

\subsection{Future work}

The most important next steps are:
\begin{enumerate}
    \item design a stronger q-ary code compactor with efficient decoding and exact failure estimates;
    \item build a sparse-LPN parameter estimator specialized to the exposed sample structure;
    \item add seeded matrix generation to shrink evaluation keys;
    \item study whether multiplication can be followed by a code-based refresh operation;
    \item explore batching or packing across rows without introducing ring/lattice assumptions;
    \item construct a public-key variant;
    \item implement a Python prototype and then a C/Rust reference implementation.
\end{enumerate}

\subsection{Recommended claim language}

For a research paper, use cautious claim language:

\begin{quote}
We propose a candidate all-code-based leveled somewhat homomorphic encryption scheme for bounded-depth arithmetic circuits.  Under sparse-LPN and q-ary LPN/code assumptions, the construction achieves compact input encryption, GSW-style nonlinear evaluation, and code/LPN-based output compaction without lattice or number-theoretic assumptions.
\end{quote}

Avoid stronger claims such as ``practical full FHE from codes'' or ``secure against all quantum attacks.''  The correct claim is conditional post-quantum plausibility under explicit code-based assumptions.


\subsection{C7 milestone status}

The C7 coordinate compactor is the conservative main candidate after the C6
relation-surface diagnostic.  It deliberately gives up the compactness of a full
projective block dictionary in exchange for removing the abundant within-block
weight-three projective relations.  This makes the current package ready as a
research prototype and manuscript milestone, while preserving a clear boundary:
it is not a production cryptosystem and does not certify concrete parameters.


\paragraph{Final v0.9-C7 status.}
The final validation pass moves the main compaction candidate from C4 projective
blocks to C7 coordinate relation-resistant compaction.  This choice sacrifices
the one-term-per-active-block advantage of C4, but removes the complete
projective weight-three relation surface discovered by C6.  The project is
therefore ready to pause as a research prototype and manuscript package, not as
a deployable cryptosystem.
